\section{Mid-Term Examination --- Set B (Units 1, 2 \& 3)}
\noindent\textbf{Course Code:} CSR101 \hfill \textbf{Max Marks:} 30\\
\textbf{Course Title:} Python Programming (AI \& Data Engg.) \hfill \textbf{Time Allowed:} 90 Minutes\\
\vspace{4mm}
\hrule
\vspace{4mm}

\subsection*{Part A: Multiple Choice Questions (15 Marks)}
\begin{enumerate}[label=\textbf{\arabic*.}]
    \item What is the output of \texttt{bool("False")}? \hfill \textbf{[Ans: \texttt{True} (non-empty string)]}
    \item Which built-in function returns an iterator yielding $(index, value)$ pairs? \hfill \textbf{[Ans: \texttt{enumerate()}]}
    \item How are keyword-only arguments denoted in a parameter list? \hfill \textbf{[Ans: After a bare \texttt{*} or \texttt{*args}]}
    \item What is the space complexity of a generator expression compared to a list comprehension? \hfill \textbf{[Ans: $\mathcal{O}(1)$ vs $\mathcal{O}(n)$]}
    \item In a class, how are class methods declared? \hfill \textbf{[Ans: With the \texttt{@classmethod} decorator and \texttt{cls} parameter]}
    \item What is the output of \texttt{[1, 2] * 2}? \hfill \textbf{[Ans: \texttt{[1, 2, 1, 2]}]}
    \item What is printed by \texttt{set([1, 2, 2, 3])}? \hfill \textbf{[Ans: \texttt{\{1, 2, 3\}}]}
    \item What is the purpose of \texttt{\_\_slots\_\_}? \hfill \textbf{[Ans: Prevents dynamic \texttt{\_\_dict\_\_} creation, saving memory]}
    \item What does \texttt{re.split(r'\textbackslash s+', text)} do? \hfill \textbf{[Ans: Splits string on one or more whitespace characters]}
    \item Can a tuple containing a list be used as a dictionary key? \hfill \textbf{[Ans: No, raises \texttt{TypeError: unhashable type: 'list'}]}
    \item What is the output of \texttt{10 // -3}? \hfill \textbf{[Ans: \texttt{-4}]}
    \item What does \texttt{zip([1, 2], ['a', 'b'])} return? \hfill \textbf{[Ans: Iterator yielding tuples: \texttt{(1, 'a'), (2, 'b')}]}
    \item Which dunder method is called when an object is used in a \texttt{with} statement? \hfill \textbf{[Ans: \texttt{\_\_enter\_\_()}]}
    \item What does \texttt{isinstance(True, int)} evaluate to? \hfill \textbf{[Ans: \texttt{True} (\texttt{bool} subclasses \texttt{int})]}
    \item What does \texttt{re.sub(r'cat', 'dog', text)} do? \hfill \textbf{[Ans: Replaces all occurrences of 'cat' with 'dog']}
\end{enumerate}

\subsection*{Part B: Subjective Coding \& Design Problems (15 Marks)}

\noindent\textbf{Question 16 [Unit 3, 5 Marks]:} Implement a custom iterator class \texttt{PrimeIterator(limit)} that yields consecutive prime numbers up to \texttt{limit} implementing the \texttt{\_\_iter\_\_} and \texttt{\_\_next\_\_} protocol.
\begin{solutionbox}
\begin{lstlisting}
class PrimeIterator:
    def __init__(self, limit: int):
        self.limit = limit
        self.current = 2

    def __iter__(self):
        return self

    def _is_prime(self, n: int) -> bool:
        for d in range(2, int(n**0.5) + 1):
            if n % d == 0: return False
        return True

    def __next__(self):
        while self.current <= self.limit:
            n = self.current
            self.current += 1
            if self._is_prime(n):
                return n
        raise StopIteration
\end{lstlisting}
\end{solutionbox}
\begin{markingrubric}
\textbf{Marking Scheme:} Class initialization and \texttt{\_\_iter\_\_} returning self (1M); Prime checking logic (2M); \texttt{\_\_next\_\_} state mutation and \texttt{StopIteration} raising (2M).
\end{markingrubric}

\noindent\textbf{Question 17 [Unit 3, 5 Marks]:} Implement a custom Vector class in 3D space supporting vector addition (\texttt{+}), dot product (\texttt{@}), and magnitude representation.
\begin{solutionbox}
\begin{lstlisting}
import math

class Vector3D:
    def __init__(self, x: float, y: float, z: float):
        self.x, self.y, self.z = x, y, z

    def __add__(self, other):
        return Vector3D(self.x + other.x, self.y + other.y, self.z + other.z)

    def __matmul__(self, other):
        return self.x * other.x + self.y * other.y + self.z * other.z

    def magnitude(self) -> float:
        return math.sqrt(self.x**2 + self.y**2 + self.z**2)

    def __repr__(self):
        return f"Vector3D({self.x}, {self.y}, {self.z})"
\end{lstlisting}
\end{solutionbox}
\begin{markingrubric}
\textbf{Marking Scheme:} \texttt{\_\_add\_\_} operator overload (1.5M); \texttt{\_\_matmul\_\_} dot product (1.5M); Magnitude and representation (2M).
\end{markingrubric}

\noindent\textbf{Question 18 [Unit 2, 5 Marks]:} Write a generator function \texttt{sliding\_window(iterable, k)} that yields consecutive sliding windows of length $k$ from a sequence.
\begin{solutionbox}
\begin{lstlisting}
from collections import deque

def sliding_window(iterable, k: int):
    it = iter(iterable)
    window = deque(maxlen=k)
    for _ in range(k):
        try:
            window.append(next(it))
        except StopIteration:
            return
    yield tuple(window)
    for elem in it:
        window.append(elem)
        yield tuple(window)
\end{lstlisting}
\end{solutionbox}
\begin{markingrubric}
\textbf{Marking Scheme:} Proper use of generator \texttt{yield} (2M); Efficient $\mathcal{O}(1)$ window updates using \texttt{deque} (2M); Edge cases handling (1M).
\end{markingrubric}
